% !TEX TS-program = pdflatex
% !TeX program = pdflatex
% !TEX encoding = UTF-8
% !TeX spellcheck = en_US

\documentclass[11pt, a4paper]{article}
%\usepackage{fullpage}
\usepackage[left=1cm,right=1cm,top=1cm,bottom=2cm]{geometry}
\usepackage[fleqn]{amsmath}
\usepackage{amssymb}
%\usepackage{indentfirst}
\usepackage[T1]{fontenc}
\usepackage[utf8]{inputenc}
\usepackage[french,english]{babel}
\usepackage{txfonts} 
\usepackage[]{graphicx}
\usepackage{multirow}
\usepackage{hyperref}
\usepackage{parskip}
\usepackage{multicol}
\usepackage{wrapfig}

\usepackage{turnstile}%Induction symbole

\usepackage{tikz}
\usetikzlibrary{arrows, automata}
\usetikzlibrary{decorations.pathmorphing}

\renewcommand{\baselinestretch}{1}

\setlength{\parindent}{24pt}


\begin{document}

%\selectlanguage {french}
%\pagestyle{empty} 

\noindent
\begin{tabular}{ll}
\multirow{3}{*}{\includegraphics[width=1.5cm]{../../../extra/logo/esi.nlp.pdf}} & \'Ecole national Supérieure d'Informatique\\
& 2\textsuperscript{nd} year second cycle (2CSSID)\\
& NLP: Natural Language Processing (2022-2023)
\end{tabular}\\[.25cm]
\noindent\rule{\textwidth}{2pt}\\[-0.5cm]
\begin{center}
{\LARGE \textbf{Lab02: Natural language detector}}
\begin{flushright}
	Abdelkrime Aries
\end{flushright}
\end{center}\vspace{-0.5cm}
\noindent\rule{\textwidth}{2pt}

We want to implement a simple language detector from scratch. 
There are many ways to implement such task; for example, using neural networks.
All those solutions can be seen as language models.
We want to understand this concept by implementing a simple n-gram model.
This model will be used to implement a language detector.

\section*{1. Program description}

Our objective is to use a Bi-gram model in order to calculate the probability of an ordered list of strings.
This probability is used to decide the language by which this list of strings is written.
The simplest design is to implement two classes:
\begin{itemize}
	\item \textbf{BiGram}: This class implements the different methods used to score a language's sentence.
	\textbf{\slshape This must be implemented}.
	\item \textbf{LangDetector}: This class uses the past one in order to train many language models; each of each language.
	\textbf{\slshape This is already implemented}.
\end{itemize}
The other objective is to implement a program dependency-free; i.e. it does not depend on any external API.


\subsection*{1.1. BiGram class}

This is the core of our program.
In our first attempt, we want just to implement a Bi-gram model with Lidstone (the value is a parameter).
To this end, the model must implement these methods: fit, score and predict.
These functions use some specific data-structures to store statistics.
The class's attributes are:
\begin{itemize}
	\item A dictionary \textbf{uni\_grams} that maps each string into an integer which is its number of occurrences.
	\item A dictionary \textbf{bi\_grams} that maps each a sequence of two strings into an integer which is its number of occurrences.
	In this case, we separate the two strings with a special character (\textbf{\#}).
	In some languages, such as python, we can use a tuple.
\end{itemize}

\subsubsection*{fit}

This function is used to train our model based on some data.
The latter is a list of lists of strings which means a list of sentences, each having a list of words or characters.
The sentences are not padded.
Its function is to fill the previous structures.

\subsubsection*{score}

Since this is a Bi-gram, we need just the current string and the past one to score the conditional probability. 
Also, it has as argument a float \textbf{alpha} which is the Lidstone factor.
In this case, we will use the trained data-structures to calculate such probability (We use log probability).
Note that the size of \textbf{uni\_grams} represents the size of vocabulary.

\subsubsection*{predict}

This method's purpose is to calculate a sentence's log probability using Bi-grams.
To this end, we pass a list of strings representing a sentence (strings can be words or characters).
Also, we pass \textbf{alpha} as Lidstone factor.
Note that the sentence is not padded.

\subsection*{1.2. LangDetector class}

\textbf{\slshape This class is already implemented}. 
This class stores many language models; either word-based or character-based.
This is a short description of its methods:
\begin{itemize}
	\item \textbf{fit}: given a list of language codes and their training text URLs, it trains many language models.
	\item \textbf{predict}: given a sentence, it calculates its log probability based on each model. 
	Then, it returns the code of the most probable language.
	\item \textbf{evaluate}: it loads a list of evaluations and evaluate our task.
\end{itemize}

\section*{3. Task training and evaluation}

This is a little description of this task's evaluation.

\subsection*{3.1. Training dataset}

In our case, we used Wikipedia articles about "Algeria".
We took just the abstract.
For each language, we created a file named "\textbf{[lang-code].train}". 
We will use just these 8 languages: ar (Arabic), en (English), es (Spanish), fa (Farsi), fr (French), it (Italian), kab (Kabyle with Latin writing system) and ur (Urdu).
Each line contains a sentence.
All the sentences do not contain the character \#.

\subsection*{3.2. Test dataset}

It is just a file named "\textbf{lang.eval}".
Each line contains a sentence, followed by \#, followed by the language's code.
For each language, we afforded 10 sentences from Wikipedia articles about "Earth".
The Urdu sentences are separated in a random way since we could not know sentence boundary character.

\subsection*{3.3. Metrics}

Since it is a multi-class classification task, we will use accuracy.


\section*{4. Questions}

Answer these questions at the beginning of your code, as comments:
\begin{enumerate}
	\item Mention one advantage of character-based language models over word-based ones in this task (language detection) and vice-versa.
	
	\item How can we enhance our n-gram model in order to better represent the probabilities?
	In this case, which model will scale better (character-based or word-based) and why?
	
	\item During prediction, does sentences' lengths affect the result? How? 
	
	\item We note that the accuracy of char-based model is 0.925 and that of word-based model is 0.125 (all models have a recall and precision of 0 except Urdu). Why Char-based model gave better performance?
	
\end{enumerate}


\section*{5. Students' Grading}

\begin{itemize}
	\item Duration: 1h (You have to return your work at the end of session)
	\item Grade
	\begin{itemize}
		\item \textbf{fit grade} (6pts) = processing all data (2pts) + all uni-grams (2pts) + all bi-grams (2pts).
		\item \textbf{score grade} (4pts) = correct and less complex uni-grams/bi-grams estimation (2pts) + applying smoothing (2pt) + log probability (1pt).
		\item \textbf{predict grade} (4pts) = correctly calculate the log probability (3pts) + correctly taking smoothing in consideration (1pt).
		\item \textbf{questions grade} (4pts) = 1pt for each question.
		\item \textbf{In time grade} (2pts) = you have 5 minutes after class to send the solution. 
		Each passed two minutes after that will subtract a 0.25pt. 
		After that, you will have a 0.
	\end{itemize}
\end{itemize}

\end{document}
